\documentclass[10pt, aspectratio=169]{beamer}

\usepackage{clases}
\usepackage[linesnumbered, noend, spanish, onelanguage]{algorithm2e}
\usepackage{bm}
\newcommand{\super}[1]{^{(#1)}}
\newcommand{\framet}[1]{\Large
#1
\vspace{1em} \\
\normalsize}
\newcommand{\dsum}{\displaystyle\sum}
\DeclareMathOperator*{\argmin}{arg\,min}

%% These macros specify information about the presentation
\title{Presentación TP de Optimización No Lineal}
\subtitle{Investigación Operativa - 2C 2025}
\author{Nazareno Faillace Mullen}
\date{}


\usefonttheme{professionalfonts} % required for mathspec
% \usepackage{mathspec}
% \setsansfont[BoldFont={Fira Sans},
% Numbers={OldStyle}]{Fira Sans Light}
% \setmathsfont(Digits)[Numbers={Lining, Proportional}]{Fira
% Sans Light}

\begin{document}

%   \shorthandoff{-}
    \frame[c]{\maketitle}

%   \AtBeginSection[]{% Print an outline at the beginning of sections
%     \begin{frame}<beamer>
%       \frametitle{Outline for Section \thesection}
%       \tableofcontents[currentsection]
%     \end{frame}}


    \AtBeginSection[]{% Print an outline at the beginning of sections
        \begin{frame}
            <beamer>
            \vfill
            \centering
            \Huge{\insertsectionhead}
            \vfill
        \end{frame}}


    \section{Regresión Lineal}

    \begin{frame}
        % \fontsize{22pt}{10pt}\selectfont
        \huge
        \[\underbrace{\text{\textbf{Regresión}}}_{\substack{\text{Queremos predecir} \\ \text{un valor continuo}}}
            {\color{red}\underbrace{\text{ \textbf{Lineal}}}_{\text{con una recta}}}\]
    \end{frame}

    \begin{frame}
        \textbf{Modelo matemático:} $Y = \beta_0 + \beta_1 X_1 + \beta_{2}X_2 + \dots + \beta_{m}X_m$
        \begin{itemize}
            \item $\beta_0$ es la ordenada al origen
            \item $\beta_1,\dots, \beta_m$ son los coeficientes
            \item $X_1,\dots, X_m$ son las variables predictoras
            \item $Y$ es la variable dependiente
        \end{itemize}
        \textbf{Modelo de regresión lineal:} $y_i = \beta_0 + \sum_{j=1}^m \beta_j x_{ij} + \varepsilon_i$
        \begin{itemize}
            \item $(x_i, y_i)$ son los datos observados
            \item $\varepsilon\super{i}$ es un error aleatorio (variación de $Y$ no explicada por $X$)
            \item $\beta_0, \beta_1, \dots, \beta_m$ son \textbf{los parámetros del modelo}
        \end{itemize}
    \end{frame}

    \begin{frame}
        \begin{columns}
            \column{0.4\textwidth}
            \textbf{Residuo:} dado $\beta$,
            definimos al residuo como la diferencia entre el valor observado ($y_i$) y el valor predicho ($\hat{y}_i$):
            \[y\super{i} - \underbrace{\left(\beta_0 + \sum_{j=0}^m\beta_j x_{ij}\right)}_{\textstyle \hat{y}_i}\]

            \centering
            \visible<2>{
                \alert{$\rightarrow$ Queremos encontrar valores para $\beta_0, \beta_1, \dots, \beta_m$
                    que \underline{minimicen} los residuos}}
            \column{0.6\textwidth}
            \includegraphics[width=\textwidth]{imagenes/regresion_ej_residuos}
        \end{columns}
    \end{frame}

    \begin{frame}
        \vspace{1em}

        \normalsize
        \textbf{Objetivo:} minimizar la suma de los residuos al cuadrado:
        \[
            RSS(\beta) = \displaystyle\sum_{i=1}^n \left(y_i - \left(\beta_0 + \sum_{j=1}^m\beta_j x_{ij}\right)
            \right)^2
        \]

        \textbf{Teorema:} si $\overline{y} = 0\Rightarrow \beta_0 = 0$.
        \vspace{1em}

        Luego, si $\overline{y} = 0$, queremos minimizar:

                \[
            RSS(\beta) = \displaystyle\sum_{i=1}^n \left(y_i - \sum_{j=1}^m\beta_j x_{ij} \right)^2
        \]
    \end{frame}

    \begin{frame}
        Sea $X = (X_1 | \dots | X_m)\in \mathbb{R}^{n\times m}$.

        Bajo algunas hipótesis, se puede probar que:
        \[
            \begin{array}{rl}
                \hat{\beta} \text{ es minimizador de } RSS(\beta) & \iff \hat{\beta} \text{ es minimizador de } \lVert X\beta - {y}\rVert_2^2 \\
                & \iff X^{T}X\hat{\beta} = X^{T}{y} \\
                & \iff \hat{\beta} \text{ es un minimizador de } \frac{1}{2}\beta^T X^{T}X\beta - X^{T}{y}\beta
            \end{array}
        \]

        Entonces para hallar $\hat{\beta}$ basta minimizar $f(\beta) = \frac{1}{2}\beta^T (X^{T}X)\beta - X^{T}y\beta$
    \end{frame}

    \begin{frame}
        \Large
        ¿Qué tan bien ajusta el modelo?
        \vspace{1em}

        \normalsize
        \textbf{Error cuadrático medio:}
        \[ECM(\beta) = \frac{1}{n}\displaystyle\sum_{i=1}^n \left(y_i - \sum_{j=1}^m\beta_j x_{ij} \right)^2\]
        \vspace{1em}

        \textbf{Coeficiente de determinación:}
        \[R^2 = 1 - \frac{\sum_{i=1}^n (y_i - \hat{y}_i)^2}
        {\sum_{i=1}^n (y_i - \overline{y})^2}\]
        Mientras más cercano a $1$ esté $R^2$, mejor.
    \end{frame}

    \begin{frame}
        \Large
        Overfitting
        \vspace{1em}

        \normalsize
        Para medir si el modelo incurre en \textit{overfitting}, dividimos a los datos en conjuntos de entrenamiento
        y de testeo.

        Si el modelo ajusta bien a los datos de entrenamiento y muy mal a los datos de testeo, entonces hay overfitting.

    \end{frame}

    \begin{frame}
        \framet{Para el TP}
        \begin{itemize}
            \item en la sección ``Carga de datos'', la función \texttt{datos\_posicion} directamente procesa la matriz
            $X$ y el vector $y$. Todos los datos de $X, y$ ya están centrados y escalados, también divididos en
            entrenamiento y testeo.
            \item la función \texttt{calcular\_errores} calcula el $R^2$ para el conjunto de entrenamiento y testeo.
        \end{itemize}
    \end{frame}

    \begin{frame}
        \framet{Problemas con Regresión Lineal Ordinaria}

        \centering
        \includegraphics[scale=0.6]{imagenes/beta_ols}
    \end{frame}

    \begin{frame}
        \framet{Regresión Lasso}

        Nueva función objetivo:
        \[f_\lambda(\beta) = \underbrace{\displaystyle\sum_{i=1}^n \left(y_i - \sum_{j=1}^m\beta_j X_{ij}
        \right)^2}_{RSS(\beta)} +
        \lambda\underbrace{\lVert\beta\rVert_1}_{L_1(\beta)}\]

        Añade un parámetro de regularización $\lambda$ para controlar la magnitud de $\beta$
        \vspace{1em}

        \textbf{Problema: } $f_\lambda$ no es diferenciable en todo punto $\Rightarrow$ $\nabla f_\lambda$ no está
        definida en todo punto.
    \end{frame}

    \begin{frame}
        Como $L_1(\beta)$ es una función convexa en $\beta$, podemos considerar las sub-derivadas:
        \[
            \partial_{\beta_j} L_1(\beta) =
            \begin{cases}
                -1 & \beta_j < 0 \\
                [-1,1] & \beta_j = 0 \\
                1 & \beta_j > 0
            \end{cases}\; \rightarrow \;
            \partial_{\beta_j} L_1(\beta) =
            \begin{cases}
                -1 & \beta_j < 0 \\
                0 & \beta_j = 0 \\
                1 & \beta_j > 0
            \end{cases}
            \Rightarrow  \partial_{\beta_j} L_1(\beta) = sg(\beta_j)
        \]

        Y podemos adaptar los métodos que estudiamos, considerando el sub-gradiente de $f_\lambda$:
        \[\partial f_\lambda(\beta) = \nabla RSS(\beta) + \lambda(\partial_{\beta_1} L_1(\beta), \dots,\partial_{\beta_m} L_1(\beta)) \]
    \end{frame}

    \begin{frame}
        \centering
        \includegraphics[scale=0.6]{imagenes/beta_l1}
    \end{frame}

    \begin{frame}
        \framet{Optimización por coordenadas}

        {\fontfamily{serif}\selectfont
        \centering
        \begin{algorithm}[H]
            \DontPrintSemicolon
            \KwIn{$f$, $\beta$, $k\_max$}
            \KwOut{$\beta$: aproximación al óptimo}
            $k\leftarrow 0$\;
            $idx \leftarrow$ índices de $\beta$ (i.e. $[0,1,\dots,m-1]$)\;
            Mezclar $idx$ \texttt{(np.random.shuffle(idx))}\;
            \While{$k < k\_max$}{
                \For{$j \in idx$}{
                    $\beta_j = \argmin_{w\in\mathbb{R}} f(\beta_1,\beta_2,\dots, \beta_{j-1}, w, \beta_{j+1}, \dots, \beta_m)$ \;
                    }
                    Mezclar $idx$\;
                    $k \leftarrow k + 1$ \;
            }
            \Return $\beta$
            \caption{Minimización por coordenadas}
        \end{algorithm}}
    \end{frame}

    \begin{frame}
        \framet{Optimización por coordenadas - Regresión Lasso}

        \textbf{Teorema: } sea $f$ convexa, $x_0$ es un minimizador de $f$ si y sólo si $0\in\partial f$.
        \vspace{1em}

        Sea $1\leq j\leq m$ y sea
        \[g(w) = f_\lambda(\beta_1,\beta_2,\dots, \beta_{j-1}, w, \beta_{j+1}, \dots, \beta_m)\]
        para minimizar $g$ basta hallar dónde se anula $\partial_{\beta_j}f_\lambda$.
        \vspace{1em}
    \end{frame}

    \begin{frame}
        Probamos que:
        \[\argmin_{w\in\mathbb{R}}
        f_\lambda(\beta_1,\beta_2,\dots, \beta_{j-1}, w, \beta_{j+1}, \dots, \beta_m) = \frac{S(\rho_j, \lambda)}{z_j}\]
        donde:
        \[
            \begin{array}{l}
                \rho_j=\dsum_{i=1}^n
                \left(y_i - \dsum_{\substack{k= 1 \\ k\neq j}}^m \beta_k x_{ik} \right)x_{ij}   \\[1em]
                z_j = \dsum_{i=1}^{n} x_{ij}^2   \\[1em]
                S(\rho_j, \lambda) =
                \begin{cases}
                    \rho_j + \lambda & \text{si $\rho_j < -\lambda$} \\
                    0 & \text{si $-\lambda \leq \rho_j \leq \lambda$} \\
                    \rho_j - \lambda & \text{si $\rho_j > \lambda$} \\
                \end{cases}
            \end{array}
        \]
    \end{frame}

    \section{Support Vector Machines}

    \begin{frame}
        \framet{SVM - Clasificación}

        \centering
        \includegraphics[scale=0.6]{imagenes/svm_1}
    \end{frame}

    \begin{frame}
        \framet{SVM - Clasificación}

        \centering
        \includegraphics[scale=0.6]{imagenes/svm_2}
    \end{frame}

    \begin{frame}
        Queremos predecir la categoría de cada individuo $y_i\in\{-1, 1\}$ a partir de sus características
        $x_i=(x_{i1}, \dots, x_{im})$.

        \textbf{Idea del método:} hallar $w\in\mathbb{R}^{m}, b\in\mathbb{R}$ para determinar el hiperplano:
        \[wx + b = 0\]
        que divida a las categorías con margen máximo.
        \vspace{1em}

        Queremos que:
        \[\begin{cases}
        (w^Tx_i + b) \leq -1 & \text{si $y_i = -1$} \\
        (w^Tx_i + b) \geq 1 & \text{si $y_i = 1$}
        \end{cases} \iff y_i(w^Tx_i + b) \geq 1\]

    \end{frame}

    \begin{frame}
        Nuestra función a minimizar:
        \[f(w,b) = \underbrace{\lambda\frac{1}{2}\lVert w\rVert_2^2}_{L_2(w)} +
        \frac{1}{n}\underbrace{\sum_{i=1}^n \max(0, 1-y_i(wx_i+b))}_{H(w,b)}\]
        \vspace{1em}
    \end{frame}

    \begin{frame}
        \centering
        \includegraphics[height=0.9\textheight]{imagenes/meme-sub-deriv}
    \end{frame}

    \begin{frame}
        En particular, tenemos que:
        \[
            \begin{array}{l}
                \partial_{w_j} \max(0, 1-y_i(w^Tx_i+b)) =
                \begin{cases}
                    0 & \text{si $y_i(w^Tx_i + b) \geq 1$} \\
                    -y_{i}x_i^{(j)} & \text{c.c.}
                \end{cases} \\[1em]
            \end{array}
        \]
        donde $x_i^{(j)}$ es la $j$-ésima coordenada de $x_i$

        \textbf{Ejercicio:} calcular $\partial_b \max(0, 1-y_i(w^Tx_i+b))$
        \vspace{1em}

        Por lo tanto, podemos aplicar los métodos que estudiamos, utilizando el sub-gradiente en vez del gradiente:
        \[
            \begin{array}{l}
                \partial_w f(w,b) = \nabla L_2(w) + \frac{1}{n}(\partial_{w_1}H(w,b),\dots, \partial_{w_j}H(w,b)) \\
                \partial_b f(w,b) = \frac{1}{n}\partial_{b}H(w,b)
            \end{array}
        \]
    \end{frame}

    \begin{frame}
        \framet{Para el TP}

        \begin{itemize}
            \item implementaremos Método de Gradiente con paso fijo $\alpha$ (parámetro) y ADAM.
            \item pueden decidir cómo implementarlos: pensar en $f(w,b)$ (diferenciando a los pesos del bias) o
            considerando todo en un vector de pesos $W=(W_0,W_1,\dots, W_j)$ donde $W_0$ sería el bias:
            \[f(W) = \lambda \frac{1}{2}\lVert (W_1,\dots, W_j)\rVert_2 ^2 + \frac{1}{n}\sum_{i=1}^n
            \max(0, 1-y_i((W_1,\dots, W_j)^Tx_i + W_0))\]
        \end{itemize}
    \end{frame}

    \begin{frame}
        \framet{¿Qué tan bien ajusta el modelo?}

        \textbf{Precisión (\textit{accuracy score}): }
        \[\dfrac{\text{clasificaciones correctas}}{n}\]
    \end{frame}
\end{document}
